\section{Coupling}
\label{sec:coupling}
A phase connection compares internal directions as the field changes. Coupling also changes the response to those directions. On the complete field space $\ComparisonSpace$, the response map is
\CouplingCouplingDisplay
The label $p$ selects the component, while $j$ runs over its partners. Their pairings supply the coefficient multiplying that component, so its response depends on both magnitude and relative alignment.

\subsection{How a change reaches the response}
For a two-component example, take
\CouplingMatrixInputExampleDisplay
Each component has one partner, giving
\CouplingMatrixResponseExampleDisplay
Changing either matrix can therefore affect both responses through their shared pairing.

The first variation must retain that dependence. Writing $L_X=D\Lambda(X)$ gives
\CouplingCouplingResponseDisplay
The first term changes the component at its current coefficient; the second changes the coefficient itself. In the example, varying only the second diagonal entry of the first component yields
\CouplingMatrixVariationExampleDisplay
The second response changes although its own input component is fixed. The derivative $L_X$ records this transmission of a change through the coupling.

\subsection{Two complex response directions}
To compare a field change with its response, place them in $\mathcal V\oplus\mathcal V$ with the product metric of Section~\ref{sec:field}. The graph consists of pairs $(v,L_Xv)$. Its perpendicular directions are obtained from the metric adjoint,
\CouplingAdjointMeaningDisplay
which transfers $L_X$ between the two entries of the pairing. The graph and normal lifts are
\CouplingPairLiftDisplay
Their comparison gives
\CouplingPairOrthogonalityDisplay
so a normal lift is perpendicular to every graph lift. Appendix~\ref{app:coupling} expands the cancellation in coordinates.

Apply both lifts to the radial and transverse directions from Section~\ref{sec:phase}. Let $\PlaneInjection$ insert their linear combinations in the chosen labelled component of $\mathcal V$. Lifting its columns gives
\CouplingPairColumnsDisplay
The graph retains $\TwoColumnFrame$ in its first position and the normal lift retains it in its second. Both planes therefore retain two independent directions, even if $L_X$ has a kernel.

Their Gram matrices satisfy
\CouplingPairGramDisplay
The zero cross-pairing records orthogonality, while the positive Gram matrices give a metric on each plane. Explicitly, for a nonzero plane-coordinate column $\CouplingPlaneCoordinates$,
\CouplingGramLengthsDisplay
Each squared length contains the strictly positive contribution from $\TwoColumnFrame\CouplingPlaneCoordinates$.

Both planes inherit the radial and transverse ordering. Their metrics and orientations therefore supply a quarter-turn on each, giving two orthogonal complex directions. In oriented orthonormal coordinates,
\FieldCouplingComplexPairDisplay
their combined length is
\FieldCouplingComplexNormDisplay
Transformations preserving this complex inner product form $\FieldCouplingUnitaryGroup$.

Preserving the complex volume form imposes the determinant-one condition. For two complex columns,
\CouplingComplexVolumeDisplay
so the unitary maps preserving this form constitute $\PairGroup$. They have matrices
\CouplingPairGroupDisplay
where $\CouplingAlpha,\CouplingBeta$ are complex and the bars denote conjugation. Metric preservation gives orthonormal columns, and determinant one retains their complex volume.

For example,
\CouplingUnitaryMixingExampleDisplay
mixes the two directions equally while retaining unit length and determinant one. Such pair transformations need not commute. Two of them are
\FieldCouplingMixingMatricesDisplay
Applied to
\FieldCouplingInitialCoordinateDisplay
in opposite orders, they give
\FieldCouplingFirstOrderDisplay
and
\FieldCouplingSecondOrderDisplay
The different results exhibit the order dependence that the shared pair actions will carry into the three-direction construction.

\subsection{The shared intermediate response}
Applying the same coupling to $Y$ produces $Z$. The derivative of this next response must pass through the change of $Y$,
\CouplingSuccessiveResponseDisplay
Here $L_Y$ is evaluated at the intermediate field. In the matrix example the next pairing is eight, giving
\CouplingSecondMatrixResponseDisplay
Thus all three levels are determined by the original matrices and the same rule; the two adjacent comparisons share $Y$.

\begin{figure}[!htbp]
\centering
\begin{minipage}[t]{.495\linewidth}\centering
\resizebox{\linewidth}{!}{\CarrierScene{2}{1}{0}}\par\smallskip
\small (a) Two response components and their pair action
\end{minipage}\hfill
\begin{minipage}[t]{.495\linewidth}\centering
\resizebox{\linewidth}{!}{\CarrierScene{3}{2}{0}}\par\smallskip
\small (b) Three coupled response components
\end{minipage}
\caption{Successive responses $X\to Y\to Z$ under $\CouplingMap$ give two adjacent pair actions, shown by solid edges. Dashed edges mark the first--third relation generated from them in Section~\ref{sec:color}.}
\label{fig:couplings}
\end{figure}
\FloatBarrier

The shared intermediate response links the two solid pair relations in Figure~\ref{fig:couplings}. Determining what their combined action generates is the next step.
\FloatBarrier
